\documentclass[10pt,a4paper]{article}
\usepackage{amsmath,amssymb,amsthm,graphicx}
\usepackage{amsmath,amssymb,amsthm,hyperref,xfrac}
\providecommand{\eg}{e.g.\ }
\providecommand{\ie}{i.e.\ }
\usepackage[scr=rsfs,cal=euler]{mathalfa}
\multlinegap0pt
\hyphenation{accu-rate applied assump-tion assump-tions decom-pos-ing defined esti-mate esti-mates exten-ded hypothe-ses para-me-ter para-me-ters remain-der required}
\newenvironment{enumeratei}
{\bgroup\def\theenumi{\roman{enumi}}\def\theenumii{\arabic{enumii}}\begin{enumerate}}
{\end{enumerate}\egroup}
\newenvironment{enumeratea}
{\bgroup\def\theenumi{\alph{enumi}}\begin{enumerate}}
{\end{enumerate}\egroup}

\newcommand{\Psfrac}[2]{(\sfrac{#1}{#2})}
\newcommand{\psfrac}[2]{\sfrac{(#1)}{#2}}
\newcommand{\spfrac}[2]{\sfrac{#1}{(#2)}}
\newcommand{\pspfrac}[2]{\sfrac{(#1)}{(#2)}}
\newcommand{\shalf}{\sfrac{1}{2}}

\newcommand{\bnorm}[1]{\bigl\|#1\bigr\|}
\newcommand{\Bnorm}[1]{\Bigl\|#1\Bigr\|}
\newcommand{\babs}[1]{\bigl|#1\bigr|}
\newcommand{\Babs}[1]{\Bigl|#1\Bigr|}

\AtBeginDocument{%
\DeclareMathOperator{\sinn}{sin}\let\sin\sinn
\DeclareMathOperator{\coos}{cos}\let\cos\coos
}

%\usepackage{enumitem}
\usepackage{physics}
%\usepackage[foot]{amsaddr}
\usepackage{bm}
\usepackage{tikz}
\usepackage{mathtools}
\usepackage{scalerel,stackengine}
\stackMath
\usetikzlibrary{arrows.meta}
\usepackage{subcaption}
%\usepackage[style=numeric]{biblatex}
%\addbibresource{main.bib}
%\setlength{\bibitemsep}{0.3\baselineskip}

\newenvironment{mycases}
{\begin{dcases*}}
{\end{dcases*}}
\theoremstyle{plain}
\newtheorem{theorem}{Theorem}[section]
\newtheorem{lemma}[theorem]{Lemma}
\newtheorem{definition}[theorem]{Definition}
\newtheorem{proposition}[theorem]{Proposition}
\newtheorem{corollary}[theorem]{Corollary}
\newtheorem*{corollary*}{Corollary}
\theoremstyle{remark}
\newtheorem{remark}[theorem]{Remark}
\counterwithout{figure}{section}
\counterwithout{equation}{section}

\newcommand{\br}[1]{\left(#1\right)}
\newcommand{\bbr}[1]{\bigl(#1 \bigr)}
\newcommand{\Bbr}[1]{\Bigl(#1 \Bigr)}
\newcommand{\bbbr}[1]{\biggl(#1 \biggr)}
\newcommand{\tpd}{\frac{1}{(2\pi)^{d}}}
\newcommand\pig[1]{\scalerel*[2.35pt]{\big#1}{%
\ensurestackMath{\addstackgap[1.2pt]{\big#1}}}}
\newcommand\pigl[1]{\mathopen{\pig{#1}}}
\newcommand\pigr[1]{\mathclose{\pig{#1}}}
\newcommand{\ak}{\abs{k}}
\newcommand{\ap}{\abs{p}}
\newcommand{\aal}{\abs{\alpha}}
\newcommand{\sbbr}[1]{\bigl[#1\bigr]}
\newcommand{\sBbr}[1]{\Bigl[#1\Bigr]}
\newcommand{\sbbbr}[1]{\biggl[#1\biggr]}
\newcommand{\sbr}[1]{\left[#1\right]}
\newcommand{\half}{\frac{1}{2}}
\newcommand{\p}{\partial}
\newcommand{\jb}[1]{\langle#1\rangle}
\newcommand{\bset}[1]{\bigl\{#1\bigr\}}
\newcommand{\Bset}[1]{\Bigl\{#1\Bigr\}}
\newcommand{\set}[1]{\left\{#1\right\}}
\newcommand{\mytextfrac}[2]{{\textstyle\frac{#1}{#2}}}
\newcommand{\qaq}{\quad\text{and}\quad}
\newcommand{\para}{\parallel}
\newcommand{\ep}{\varepsilon}
\newcommand{\R}{\mathbb{R}}
\newcommand{\C}{\mathbb{C}}
\newcommand{\al}{\alpha}
\newcommand{\intRd}{\int_{\R^d}}
\newcommand{\intRdk}{\int_{\R^d_k}}
\newcommand{\intRdp}{\int_{\R^d_p}}
\newcommand{\lam}{\lambda}
\newcommand{\ceils}[1]{\lceil#1\rceil}
\newcommand{\Qin}{Q_{\mathrm{in}}}
\newcommand{\hin}{h_{\mathrm{in}}}
\newcommand\reallywidehat[1]{%
\savestack{\tmpbox}{\stretchto{%
\scaleto{%
\scalerel*[\widthof{\ensuremath{#1}}]{\kern.pt\mathchar"0362\kern.1pt}%
{\rule{0ex}{\textheight}}%WIDTH-LIMITED CIRCUMFLEX
}{\textheight}%
}{2.4ex}}%
\stackon[-6.9pt]{#1}{\tmpbox}%
}
\title{Derivative bounds in the velocity variable for the linear and nonlinear Wigner terms}
\author{Marnie Smith}
\date{doi:10.5802/jep.353; CC BY4.0}
\begin{document}
\maketitle
\noindent\textit{Editorial note.} The counted statement is the author's labelled lemma `etalemma` (source lines 715-728); the proof environment begins on the next line after \textbackslash{}end\{lemma\}, so the pairing is by adjacency.\par\medskip
\noindent\textit{Editorial note.} The retained context is the author's own lead-in, including the display that defines the terms L\_P and N\_P and the remark about the order O(hbar) terms that the statement's second summand encodes.\par\medskip
\noindent\textit{Editorial note.} No counter adjustment is retained: the publisher numbering is not reproduced and the item prints with the wrapper's own sequential numbering.\par\medskip
\noindent\textit{Editorial note.} No retained line refers to a label, citation or figure outside the retained span, so the package declares no external reference; the retained preamble is the publisher layout adapter plus the native macro, operator and theorem block of the same pinned source.\par\medskip
\noindent\textit{Editorial scope.} One original result of the pinned publisher article, counted once: the velocity-derivative bounds for the linear and nonlinear Wigner terms (author label etalemma).\par\medskip
\noindent\textit{Allowed standard prerequisites (editorial).}\begin{itemize}
\item The linear and nonlinear Wigner terms. L\_P and N\_P are the two terms of the Wigner equation, defined by the displays preceding the lemma, with L\_P carrying the oscillatory factor sin((hbar/2) k.(eta - k t)).. Both bounds in the statement are about the derivatives in eta of the Wigner transforms of L\_P and N\_P.
\item Leibniz rule for multi-index derivatives. D\textasciicircum{}alpha of a product expands as a sum over beta <= alpha with multinomial coefficients.. It is the first step of the proof at source lines 730-737.
\item Elementary trigonometric bounds. |sin x| <= |x| and |sin x|, |cos x| <= 1.. They separate the beta = alpha term from the beta < alpha terms at source lines 738-745 and produce the two different powers of hbar |k|.
\end{itemize}\medskip

Essential to the analysis are bounds on derivatives of the Wigner transform in the \( \eta \) variable, \( D^\alpha_\eta\widehat{W[P]} \), corresponding to velocity weights \( \bm{\xi} \) on \( P \). For multi-indices $\alpha$ with $\abs{\alpha}\geq2$, additional terms of order \( O(\hbar) \) arise from differentiation of the oscillatory phase, in stark contrast to the classical case.

\begin{lemma}\label{etalemma}
Let $\alpha$ be a multi-index. The Wigner transform of ${\mathrm{L}_P}$ satisfies\vspace*{-5pt}
\begin{multline*}
\babs{D^\alpha_\eta\widehat{W[{\mathrm{L}_P}]}(t,k,\eta)}\lesssim \ak\abs{\widehat w(k)}\abs{\widehat{\rho_Q}(t,k)}\abs{\eta-kt}\babs{D^\alpha_\eta\widehat g(\eta-kt)}\\
+\ak\jb{\hbar k}^{\abs{\alpha}-1}\abs{\widehat w(k)}\abs{\widehat{\rho_Q}(t,k)}\sum_{\abs{\beta}\leq \aal-1}\abs{D^\beta_\eta\widehat g(\eta-kt)}.
\end{multline*}
The Wigner transform of ${\mathrm{N}_P}$ satisfies\vspace*{-5pt}
\begin{multline*}
\babs{D^\alpha_\eta\widehat{W[{\mathrm{N}_P}]}(t,k,\eta)}\\
\lesssim\int_\ell \abs{\ell }\,\abs{\widehat w(\ell )}\,\abs{\widehat{\rho_Q}(t,\ell )}\,\abs{\eta-kt}\,\babs{D_\eta^\alpha \widehat{W[P]}(t,k-\ell ,\eta-\ell t)}\,d\ell\\
+\sum_{\abs{\beta}\leq \abs{\alpha}-1}\int_\ell\abs{\ell}\,\jb{\hbar \ell}^{\abs{\alpha}-1}\,\abs{\widehat w(\ell )}\,\abs{\widehat{\rho_Q}(t,\ell )}\,\babs{D_\eta^\beta \widehat{W[P]}(t,k-\ell ,\eta-\ell t)}\,d\ell.
\end{multline*}
When $\abs{\alpha}=0$, sums over $\abs{\beta}\leq\abs{\alpha}-1$ are understood to be empty.
\end{lemma}

\begin{proof}
Applying the Leibniz rule to the linear term gives\vspace*{-3pt}
\[
D_\eta^\alpha\sbr{\sin\br{\mytextfrac{\hbar}{2}k\cdot(\eta-kt)}\widehat g(\eta-kt)}
=\sum_{\beta\leq \alpha}\begin{pmatrix}
\alpha\\
\beta
\end{pmatrix}D_\eta^{\alpha-\beta}\sbr{\sin\br{\mytextfrac{\hbar}{2}k\cdot(\eta-kt)}}D^\beta_\eta\widehat g(\eta-kt).
\]
For the term with \( \beta = \alpha \), use \( \abs{\sin x} \leq \abs{x} \) to get\vspace*{-3pt}
\begin{align*}
\abs{\sin\br{\mytextfrac{\hbar}{2}k\cdot(\eta-kt)}}\leq \hbar\abs{k}\abs{\eta-kt}.
\end{align*}
For terms with \( \beta < \alpha \), differentiate $\sin$ directly and use \( \abs{\sin x},\abs{\cos x} \leq 1 \) to obtain\vspace*{-3pt}
\begin{align*}
\abs{D^{\alpha-\beta}_{\eta}\sbr{\sin\br{\mytextfrac{\hbar}{2}k\cdot(\eta-kt)}}}\leq (\hbar \abs{k})^{\abs{\alpha-\beta}}.
\end{align*}
This leads to the estimate
\begin{multline*}
\abs{D_\eta^\alpha\sbr{\bbr{e^{-i\Psfrac{\hbar}{2}k\cdot(\eta-kt)}-e^{i\Psfrac{\hbar}{2}k\cdot(\eta-kt)}}\widehat g(\eta-kt)}}\\
\lesssim \hbar\ak\abs{\eta-kt}\abs{D_\eta^\alpha\widehat g(\eta-kt)}+\sum_{\beta<\alpha}(\hbar\ak)^{\abs{\alpha-\beta}}\abs{D_\eta^\beta\widehat g(\eta-kt)}.
\end{multline*}
Since for any \( \beta < \alpha \), $(\hbar |k|)^{|\alpha - \beta|} \lesssim \hbar |k| \jb{\hbar k}^{|\alpha| - 1}$, this gives the bound for $\widehat{W[{\mathrm{L}_P}]}$. The estimate for \( \widehat{W[{\mathrm{N}_P}]} \) follows similarly.\end{proof}

\begin{thebibliography}{99}
\bibitem{r1} J. Bedrossian - “Nonlinear echoes and Landau damping with insufficient regularity”, Tunis. J. Math. 3 (2021) no. 1, p. 121-205
\bibitem{r2} J. Bedrossian, N. Masmoudi \& C. Mouhot - “Landau damping: paraproducts and Gevrey regularity”, Ann. PDE 2 (2016) no. 1, article ID 4, 71 pages
\bibitem{r3} J. Bedrossian, N. Masmoudi \& C. Mouhot - “Landau damping in finite regularity for unconfined systems with screened interactions”, Comm. Pure Appl. Math. 71 (2018) no. 3, p. 537-576
\bibitem{r4} J. Bedrossian, N. Masmoudi \& C. Mouhot - “Linearized wave-damping structure of Vlasov–Poisson in $\mathbb{R}^3$”, SIAM J. Math. Anal. 54 (2022) no. 4, p. 4379-4406
\bibitem{r5} A. Borie, S. Hadama \& J. Sabin - “Scattering for the positive density Hartree equation”, 2025
\bibitem{r6} A. Bove, G. Da Prato \& G. Fano - “An existence proof for the Hartree-Fock time-dependent problem with bounded two-body interaction”, Comm. Math. Phys. 37 (1974), p. 183-191
\bibitem{r7} A. Bove, G. Da Prato \& G. Fano - “On the Hartree-Fock time-dependent problem”, Comm. Math. Phys. 49 (1976) no. 1, p. 25-33
\bibitem{r8} T. Cazenave - Semilinear Schrödinger equations, Courant Lecture Notes in Math., vol. 10, American Mathematical Society, Providence, RI, 2003
\bibitem{r9} J. M. Chadam - “The time-dependent Hartree-Fock equations with Coulomb two-body interaction”, Comm. Math. Phys. 46 (1976) no. 2, p. 99-104
\bibitem{r10} T. Chen, Y. Hong \& N. Pavlović - “Global well-posedness of the NLS system for infinitely many fermions”, Arch. Rational Mech. Anal. 224 (2017) no. 1, p. 91-123
\bibitem{r11} T. Chen, Y. Hong \& N. Pavlović - “On the scattering problem for infinitely many fermions in dimensions $d\ge 3$ at positive temperature”, Ann. Inst. H. Poincaré C Anal. Non Linéaire 35 (2018) no. 2, p. 393-416
\bibitem{r12} J. J. Chong, L. Lafleche \& C. Saffirio - “On the $L^2$ rate of convergence in the limit from the Hartree to the Vlasov-Poisson equation”, J. Éc. polytech. Math. 10 (2023), p. 703-726
\bibitem{r13} C. Collot \& A.-S. de Suzzoni - “Stability of equilibria for a Hartree equation for random fields”, J. Math. Pures Appl. (9) 137 (2020), p. 70-100
\bibitem{r14} C. Collot \& A.-S. de Suzzoni - “Stability of steady states for Hartree and Schrödinger equations for infinitely many particles”, Ann. H. Lebesgue 5 (2022), p. 429-490
\bibitem{r15} E. Grenier, T. T. Nguyen \& I. Rodnianski - “Landau damping for analytic and Gevrey data”, Math. Res. Lett. 28 (2021) no. 6, p. 1679-1702
\bibitem{r16} S. Hadama - “Asymptotic stability of a wide class of stationary solutions for the Hartree and Schrödinger equations for infinitely many particles”, Ann. H. Lebesgue 8 (2025), p. 181-218
\bibitem{r17} A. D. Ionescu, B. Pausader, X. Wang \& K. Widmayer - “Nonlinear Landau damping for the Vlasov-Poisson system in $\mathbb{R}^3$: the Poisson equilibrium”, Ann. PDE 10 (2024) no. 1, article ID 2, 78 pages
\bibitem{r18} L. Landau - “On the vibrations of the electronic plasma”, Acad. Sci. USSR. J. Phys. 10 (1946), p. 25-34
\bibitem{r19} M. Lewin \& J. Sabin - “The Hartree equation for infinitely many particles, II: Dispersion and scattering in 2D”, Anal. PDE 7 (2014) no. 6, p. 1339-1363
\bibitem{r20} M. Lewin \& J. Sabin - “The Hartree equation for infinitely many particles I. Well-posedness theory”, Comm. Math. Phys. 334 (2015) no. 1, p. 117-170
\bibitem{r21} M. Lewin \& J. Sabin - “The Hartree and Vlasov equations at positive density”, Comm. Partial Differential Equations 45 (2020) no. 12, p. 1702-1754
\bibitem{r22} C. Mouhot \& C. Villani - “On Landau damping”, Acta Math. 207 (2011) no. 1, p. 29-201
\bibitem{r23} T. T. Nguyen \& C. You - “Modified scattering for long-range Hartree equations of infinite rank near vacuum”, SIAM J. Math. Anal. 57 (2025) no. 6, p. 6486-6497
\bibitem{r24} T. T. Nguyen \& C. You - “Plasmons for the Hartree equations with Coulomb interaction”, Probab. Math. Phys. 6 (2025) no. 3, p. 913-960
\bibitem{r25} E. M. Stein \& G. Weiss - “Interpolation of operators with change of measures”, Trans. Amer. Math. Soc. 87 (1958), p. 159-172
\bibitem{r26} C. You - “Phase mixing estimates for the nonlinear Hartree equation of infinite rank”, 2024
\end{thebibliography}
\end{document}
